%%%%%%%%%%  scan idx 259  |  folio 246  %%%%%%%%%%
\SGAPrintPage{246}%
\sgaStableTarget{sga6.e03.prop.3-6}\textbf{Proposition 3.6.} \uline{Let}
\[
\begin{tikzcd}
X \arrow[r,hook,"i"] \arrow[d,"f"'] & X' \arrow[dl,"f'"]\\
Y &
\end{tikzcd}
\]
\uline{be a commutative triangle of schemes}, \uline{where} \(i\) \uline{is a closed immersion and} \(f'\) \uline{a smooth morphism}, \uline{and let} \(E\in\Ob D(X)\). \uline{For} \(E\) \uline{to be locally of finite flat amplitude rel.} \uline{to} \(f\), \uline{it is necessary and sufficient that} \(i_*(E)\) \uline{be locally of finite flat amplitude} (\uline{in the absolute sense}). \uline{More precisely}, \uline{if} \(f'\) \uline{is of relative dimension} \(\leq d\), \uline{we have}
\begin{align*}
\text{(i)}\quad &\tamp(i_*(E))\subset[a,b] \Longrightarrow \tamp_f(E)\subset[a,b],\\
\text{(ii)}\quad &\tamp_f(E)\subset[a,b] \Longrightarrow \tamp(i_*(E))\subset[a-d,b].
\end{align*}
%%%%%%%%%%  scan idx 260  |  folio 247  %%%%%%%%%%
\SGAPrintPage{247}%
\uline{Proof}. For an interval \([m,n]\) of \(\mathbb Z\), we evidently have
\[
 \tamp_f(E)\subset[m,n]
 \Longleftrightarrow
 \tamp_{f'}(i_*(E))\subset[m,n],
\]
so that we may suppose \(X=X'\), \(i=\Id_X\), \(f=f'\). Relation (i) follows trivially from (3.4.1 a)). Let us prove (ii). We have a Cartesian square
\[
\begin{tikzcd}
X \arrow[d,"f"'] & X\times_Y X \arrow[l,"\operatorname{pr}_2"'] \arrow[d,"\operatorname{pr}_1"]\\
Y & X \arrow[l,"f"']
\end{tikzcd},
\]
and the diagonal morphism \(s:X\longrightarrow X\times_YX\) is a section of \(\operatorname{pr}_1\) and \(\operatorname{pr}_2\). Suppose \(f\) is smooth of relative dimension \(d\). Then (VII 1.10) \(s\) is a regular immersion of codimension \(\leq d\), hence is of tor-dimension \(\leq d\) (as is seen using the Koszul complex). By (3.5.2), we have
\[
 \tamp_f(E)\subset[a,b]
 \Longrightarrow
 \tamp_{\operatorname{pr}_1}(\operatorname{pr}_2^*(E))\subset[a,b],
\]
and by (3.4.1 b)) we have
\[
 \tamp_{\operatorname{pr}_1}(\operatorname{pr}_2^*(E))\subset[a,b]
 \Longrightarrow
 \tamp(\mathrm Ls^*\operatorname{pr}_2^*(E))\subset[a-d,b],
\]
giving implication (ii), since \(\mathrm Ls^*\operatorname{pr}_2^*\simeq\Id\).

\sgaStableTarget{sga6.e03.rem.3-6-1}\textbf{Remark 3.6.1.} Statement (3.6) would be false with only the assumption that \(f'\) is flat; see (3.4.2).

\sgaStableTarget{sga6.e03.prop.3-7}\textbf{Proposition 3.7.} (``\uline{Projection formula}'')\SGAAuthorFootnote{1}{cf. (H II 5.6).}. \uline{Let} \(S\) \uline{be a quasi-compact scheme}, \(f:X\to Y\) \uline{a quasi-compact and quasi-separated} \(S\)-\uline{morphism of quasi-compact schemes over} \(S\), \(E\in\Ob D^b(X)_{\qcoh}\), \(M\in\Ob D^b(S)_{\qcoh}\)\SGAAuthorFootnote{2}{The subscript ``qcoh'' indicates the full subcategory formed by complexes with quasi-coherent cohomology.}. \uline{If} \(M\) \uline{is} \uline{of finite flat amplitude}, \uline{or if} \(E\) \uline{is of finite flat amplitude rel.} \uline{to} \(S\), \uline{there}
%%%%%%%%%%  scan idx 261  |  folio 248  %%%%%%%%%%
\SGAPrintPage{248}%
\uline{is a canonical functorial isomorphism in} \(D(Y)_{\qcoh}\)
\[
 Rf_*(E)\otimes_S^{\mathrm L}M
 \xrightarrow{\sim}
 Rf_*(E\otimes_S^{\mathrm L}M)\qquad(\text{cf. }(1.6)).
\]

\uline{Proof}. First note that there is (EGA III 1.4.12 and IV 1.7.21) an integer \(N\) such that \(R^if_*(F)=0\) for every \(i>N\) and every quasi-coherent \(\underline{O}_X\)-Module \(F\). Consequently the functor \(Rf_*\) sends \(D^b(X)_{\qcoh}\) into \(D^b(Y)\) (and even \(D^b(Y)_{\qcoh}\) (EGA III 1.4.10 and IV 1.7.21)). Suppose \(E\) (resp. \(M\)) is of finite flat amplitude rel. to \(S\) (resp. finite flat amplitude). The adjunction arrow \(Lf^*Rf_*(E)\longrightarrow E\) defines, by tensoring with \(M\), a canonical arrow of \(D^-(X)\)
\[
 Lf^*(Rf_*(E)\otimes_S^{\mathrm L}M)
 \longrightarrow E\otimes_S^{\mathrm L}M,
\]
where \(E\otimes_S^{\mathrm L}M\in\Ob D^b(X)\). By adjunction we deduce a canonical arrow of \(D^-(Y)\)
\sgaStableTarget{sga6.e03.eq.unnamed}\[
 Rf_*(E)\otimes_S^{\mathrm L}M
 \longrightarrow Rf_*(E\otimes_S^{\mathrm L}M).\tag{*}
\]
Its target and source moreover have quasi-coherent cohomology, since \(Rf_*\) sends \(D^b(X)_{\qcoh}\) into \(D^b(Y)_{\qcoh}\) (EGA III 1.4.10 and IV 1.7.21), and \(\otimes_S^{\mathrm L}\) sends \(D^-(X)_{\qcoh}\times D^-(S)_{\qcoh}\) into \(D^-(X)_{\qcoh}\) (likewise with \(Y\) in place of \(X\)). Let us prove that \((*)\) is an isomorphism. The question being local on \(Y\), we may suppose \(S\) and \(Y\) affine. By dévissage, in the first case we immediately reduce to the case where \(M\) is an ordinary quasi-coherent \(\underline{O}_S\)-Module. Noting that the functors \(Rf_*(E)\otimes_S^{\mathrm L}.\) and \(Rf_*(E\otimes_S^{\mathrm L}.)\) are ``way-out to the left'' on the category \(D^b(S)_{\qcoh}\) ([H] I 7), we then reduce, by taking a left resolution of \(M\) by free Modules and truncating, to the case where \(M\) is free. Since the \(R^if_*\) commute with inductive limits (\(f\) being quasi-compact and quasi-separated), we may even suppose \(M\) free of finite type, and then \(M=\underline{O}_S\). The arrow \((*)\) is then tautologically an isomorphism. In the second case, by dévissage we reduce to \(M=P\), where \(P\) is a flat \(A\)-module
%%%%%%%%%%  scan idx 262  |  folio 249  %%%%%%%%%%
\SGAPrintPage{249}%
\((S=\Spec(A))\). Putting \(A'=D_A(P)\) (EGA \(0_{\mathrm{IV}}\) 18.2.3), \(S'=\Spec(A')\), everything reduces to proving that \(Rf_*(E)\) commutes with the flat base change \(S'\to S\), i.e. that the canonical arrow
\[
 q^*Rf_*(E)\longrightarrow Rf'_*p^*(E)
\]
is an isomorphism (we have put \(f_{(S')}=f'\) and denoted by \(p\) (resp. \(q\)) the canonical projection \(X_{(S')}\longrightarrow X\) (resp. \(Y_{(S')}\longrightarrow Y\))). But this follows from (EGA III 1.4.15 and IV 1.7.21). This completes the proof of (3.7).

\sgaStableTarget{sga6.e03.cor.3-7-1}\textbf{Corollary 3.7.1.} \uline{Let} \(S\) \uline{be a scheme}, \(f:X\to Y\) \uline{a quasi-compact and quasi-separated} \(S\)-\uline{morphism of quasi-compact schemes over} \(S\), \(E\in\Ob D^b(X)_{\qcoh}\). \uline{If} \(E\) \uline{is locally of finite flat amplitude rel.} \uline{to} \(S\), \uline{the same holds for} \(Rf_*(E)\). \uline{More precisely}, \uline{if} \(S\) \uline{is quasi-compact}, \uline{and if} \(N\) \uline{is an integer such that} \(R^if_*(F)=0\) \uline{for every} \(i>N\) \uline{and every quasi-coherent} \(\underline{O}_X\)-\uline{Module} \(F\) (EGA III 1.4.12 \uline{and} IV 1.7.21), \uline{then we have}
\[
 \tamp_S(E)\subset[a,b]
 \Longrightarrow
 \tamp_S(Rf_*(E))
 \subset[a,b+N].
\]

\uline{Proof}. We may evidently restrict ourselves to proving the last assertion. By (3.3 (i)\(\Longleftrightarrow\)(iv)), and taking into account that \(Rf_*\) ``commutes with restriction to an open subset of \(S\)'', we must prove that the hypothesis \(\tamp_S(E)\subset[a,b]\) implies
\[
 H^i(Rf_*(E)\otimes_S^{\mathrm L}M)=0
\]
for every \(i\notin[a,b+N]\) and every quasi-coherent \(\underline{O}_S\)-Module \(M\). Now the preceding hypothesis means that \(H^i(E\otimes_S^{\mathrm L}M)=0\) for every \(i\notin[a,b]\) and every quasi-coherent \(\underline{O}_S\)-Module \(M\), and (3.7) gives the result.
%%%%%%%%%%  scan idx 263  |  folio 250  %%%%%%%%%%
\SGAPrintPage{250}%
\sgaStableTarget{sga6.e03.cor.3-7-2}\textbf{Corollary 3.7.2.} \uline{Let} \(f:X\to Y\) \uline{be a quasi-compact}, \uline{quasi-separated morphism}, \uline{locally of finite tor-dimension} (3.1), \uline{and let} \(E\in\Ob D^b(X)_{\qcoh}\). \uline{If} \(E\) \uline{is locally of finite flat amplitude} (\uline{in the absolute sense}), \uline{the same holds for} \(Rf_*(E)\).

\uline{Proof}. Combine (3.4.1 a)) and (3.7.1).

\sgaStableTarget{sga6.e03.sec.4}\section*{4. Relative perfection}

\sgaStableTarget{sga6.e03.def.4-1}\textbf{Definition 4.1.} \uline{Let} \(f:X\to Y\) \uline{be a locally finite type morphism of schemes}, \(E\in\Ob D(X)\), \([a,b]\) \uline{an interval of} \(\mathbb Z\). \uline{We say that} \(E\) \uline{is} of perfect amplitude \uline{relative to} \(f\) (\uline{or} \(Y\) \uline{when there is no ambiguity}), contained in \([a,b]\), \uline{and write}
\[
 \parfamp_f(E)\subset[a,b]\qquad(\text{or }\parfamp_Y(E)\subset[a,b])
\]
\uline{if}, \uline{relative to} \(f\), \(E\) \uline{is pseudo-coherent} (1.2) \uline{and of flat amplitude contained in} \([a,b]\) (3.1). \uline{We say that} \(f\) is of perfect amplitude contained in \([-n,0]\), \uline{or of} perfect dimension \(\leq n\) \((n\geq0)\), \uline{if} \(\underline{O}_X\) \uline{is of perfect amplitude contained in} \([-n,0]\). \uline{We say that} \(E\) \uline{is} of finite perfect amplitude relative to \(f\) \uline{if there is an interval} \([a,b]\) \uline{such that} \(\parfamp_f(E)\subset[a,b]\). \uline{We say that} \(E\) \uline{is} perfect rel. to \(f\) \uline{if} \(E\) \uline{is locally of finite perfect amplitude rel.} \uline{to} \(f\). \uline{We say that} \(f\) is perfect \uline{if} \(f\) \uline{is} \uline{locally of finite perfect amplitude}.

\sgaStableTarget{sga6.e03.ex.4-1-1}\textbf{Example 4.1.1.} A regular immersion of codimension \(d\) (VII 1.4) is of perfect dimension \(\leq d\) (as is seen from the Koszul complex). A locally complete-intersection morphism (i.e. locally a composite
